\documentclass{article}
\usepackage{nips15submit_e,times}
\usepackage{hyperref}
\usepackage{url}

\usepackage[table]{xcolor}
\usepackage{colortbl}


\title{Learning to Teach}


\author{
Gordon Chen\\
School of Computing\\
University of Connecticut\\
\texttt{gbc24001@uconn.edu}\\
}

\newcommand{\fix}{\marginpar{FIX}}
\newcommand{\new}{\marginpar{NEW}}

\nipsfinalcopy

\begin{document}


\maketitle

\begin{abstract}
Despite recent improvements in LLM capabilities, LLMs today can often be poor teachers.
They give explanations that are too difficult to understand or not detailed enough, and give away
solutions too quickly.
In this project, we experiment with using reinforcement learning to fine-tune a teacher LLM to give
feedback that improves a frozen student LLM's performance through in-context learning.
When a student gets a problem wrong, the teacher sees the question and the students attempt,
and produces feedback. That feedback is added to the student's context, and we evaluate the student
on questions that the teacher has never seen but test the same concept.
The teacher is then rewarded if the student's performance increases relative to a baseline performance
without teacher's feedback, which pushes the teacher towards guidance that helps with generalization
instead of overfitting to a single problem.
\end{abstract}

\section{Problem}
LLMs today are not native teachers. They give out answers too quickly, give explanations that are
miscalibrated for the user's knowledge level, and in general serve the role of an encyclopedia more
than that of a teacher who knows how to help a student along.

While it's not fair to blame LLMs for lazy students just getting the answer from LLMs, it is a problem
that model developers can address when students that want to learn are deprived of the ability
to struggle through a problem when an LLM gives them the solution immediately instead of nudging them
in the right direction.

\section{Goal}
The goal of this project is to train an LLM to be the best possible teacher, to natively optimize
for teaching ability so that students can benefit the most from trying to learn from LLMs.

Essentially I want use reinforcement learning to fine-tune a teacher LLM to provide the most useful
feedback for a student LLM. The student LLM would be frozen, but would be able to learn in-context from
the teacher's feedback. For every problem a student LLM gets wrong, the teacher LLM would be given
the question and the student LLM's response, and ouput feedback for the student. That feedback is then placed
in the student's context, and the student is evaulated on similar questions generated by a frozen,
question-writing LLM. As a baseline, the student would also be evaulated on those questions without
the teacher's feedback in its context, and the difference in the student's accuracy would serve as a learning
signal for the teacher to know if its feedback is helpful or not.

By evaluating the student on similar problems the teacher has never seen, we prevent the teacher from just
giving the student the answer and calling it a day because the teacher's answer must help the student
generalize to the similar questions.

If all goes well and I have time, I would like to explore multi-turn interactions between
the teacher and the student. Instead of just giving feedback a single time, the student would generate
a solution, the teacher would give feedback, the student would do a similar problem with the feedback,
and the whole process would repeat. I think that this would be a more valuable task to train on because
in the real world, teachers help students multiple times, correcting their mistakes as they come up in
different questions.

\section{Importance}
This is an important project because as LLMs are getting smarter, their potential to become
really cheap, widely available, infinitely-patient, competent personal tutors is increasing.
However, currently many students are using LLMs in a very destructive way, using LLMs to offload
the need to think and learn by just asking LLMs to do their homework. For all of the
work going into making LLMs smarter, in my opinion, not enough focus has been put into making the information
transfer between LLMs (which are getting smarter very quickly) and humans more efficient. I hope that this project
shows that reinforcement learning can teach LLMs to be really good teachers, and that there is a lot of
low hanging fruit in terms of making them more effective teachers.

\section{Expected Results}
I hope to see that fine-tuning the teacher with my proposed pipeline improves the student's ability
to answer questions correctly. This would show that the teacher is learning to teach the student better
and better, learning to identify what the student is doing wrong and telling them how to correct their behavior.
I would be able to see this by there being a large difference between the student's accuracy with and without
the teacher's feedback.

I also hope to be able to identify higher-level teaching methods that the teacher LLM learns to use.
For example, perhaps the Socratic style of asking the student questions might emerge from training,
or perhaps the teacher might learn some other more effective way to elicit correct problem-solving behavior
from students. If I could identify the teaching methods that the teacher LLM learns are most useful,
then in my own prompting I could use those teaching methods to more effectively teach LLMs skills in-context.
I also wonder if the methods the teacher LLM uses can be applied to teaching humans as well.

\section{Data}
I plan on synthetically generating a dataset with multiple concepts and multiple questions per concept.
While this paradigm of learning to teach could potentially work on any subject where it is possible to
objectively know if a student is doing better (multiple choice is possible for English and general knowledge),
I plan on training a teacher on math. The dataset would then have multiple relatively simple math techniques,
such as solving quadratic equations, math word problems, and solving systems of equations, that I can hopefully
easily synthetically generate, create similar but diverse enough variations of, and be relatively easy to verify.

Since the goal of this project is to see if a teacher can be trained to improve student performance,
these synthetically generated problems need to be hard enough that student performance is initally bad,
but easy enough that with the proper teacher guidance, it is possible for a small student model to find
the solution.

\section{Algorithms}
I plan to use GRPO and LoRA to fine-tune the teacher LLM. By using LoRA and fine-tuning a model instead
of training a model from scratch, I hope to make this project feasible to develop on my RTX 3060, or
at least not cost too many cloud compute credits.

While I am only fine-tuning the teacher LLM, I also need to be able to sample from a (potentially smaller)
student LLM, as well as a question writing model, although perhaps some problems can be programatically
generated with substitution instead of being written by a question writing LLM.

\section{Schedule and Milestones}
\renewcommand{\arraystretch}{1.2}
\begin{tabular}{ |l|l|l| } 
 \hline
 3/6 & \cellcolor{red} Proposal due & \\
 \hline

 Week of 3/13 & & Background research: GRPO, LoRA. Generate / find dataset.\\
 \hline
 Week of 3/20 & \cellcolor{green} Spring Break & Start trying to train teacher, write report.\\
 \hline

 3/20 & \cellcolor{red} Midpoint check & 4-page report: progress + issues\\
 \hline

 Week of 3/27 & & Continue trying to train teacher.\\
 \hline
 Week of 4/3 & & $\downarrow$\\
 \hline
 Week of 4/10 & & $\downarrow$\\
 \hline
 Week of 4/17 & & Multi-turn experiments?\\
 \hline
 Week of 4/24 & & Work on presentation\\
 \hline

 4/23 - 4/30 & \cellcolor{red} Presentation & 15-min presentation: problem formulation, algorithms, results, demo \\
 \hline

 Week of 5/1 & & Work on final report\\
 \hline
 Week of 5/8 & & $\downarrow$\\
 \hline

 5/8 & \cellcolor{red} Final report & 6-page report: problem formulation, algs, results. \\
 \hline
\end{tabular}

% \subsubsection*{Acknowledgments}
% \subsubsection*{References}
\end{document}
